% paper/sections/01_introduction.tex

\section{はじめに：問題の背景・数値的困難・本稿の構成}
\label{sec:intro}

% ─────────────────────────────────────────────────────────────────────────────
% 【論文の芯】本稿を貫く単一の設計哲学
% ─────────────────────────────────────────────────────────────────────────────

気液二相流の数値シミュレーションでは，界面位置，物性ジャンプ，表面張力，
圧力投影を同じ離散構造の上で扱えなければならない．
この整合が崩れると，静止すべき液滴界面に人工的な渦流れ
（\textbf{寄生流れ，spurious currents}）が現れる~\cite{Popinet2009, Francois2006}．
CSF（連続表面力）モデル~\cite{Brackbill1992} を用いる場合には，
界面幅パラメータ $\varepsilon$
（本稿の CLS スケーリングでは $\varepsilon = C_\varepsilon\,\Delta x_{\min}$，
$C_\varepsilon \approx 1.0$--$2.0$；§~\ref{sec:reinit}）に由来する
$\Ord{\varepsilon^2} \approx \Ord{h^2}$ の拡散界面モデル誤差も残る．
したがって，本稿が対象とする問題は単に高次差分を導入することではなく，
表面張力・圧力・速度補正・界面移流を同じ面上の離散構造で閉じ，
どの誤差が離散化で除去され，どの誤差がモデルまたは検証範囲として残るかを明確に分離することである．

\noindent\textbf{本稿の設計原理．}
本稿の中核は，界面幾何を一度作って終わりにせず，
圧力勾配，表面張力，速度補正，界面移流が同じ面値・同じ面係数・同じ発散演算子を参照するように
閉じることである．
この組を以下では\emph{面上の共通離散構造}と呼ぶ．
また，圧力射影後に保存形 CLS 移流へ渡す標準の面速度を
\emph{射影後の面速度}と呼ぶ．
節点値として持つべき量と，面上の仕事・加速度として釣り合わせるべき量を分けることも
本稿の重要な約束である．
導入部ではこの区別を「面で同じ量を釣り合わせる」という物理的な言葉で述べ，
離散微分形式としての厳密な記法は，圧力・速度連成を扱う第~\ref{sec:collocate}章以降で導入する．

この原理に基づき，本稿では二つの圧力・表面張力経路を明確に分ける．
第一は，CSF 体積力を用いる一流体の対照経路である．
ここでは Balanced--Force 条件
（本稿で系統化する 7 原則 P-1--P-7；§~\ref{sec:bf_seven_principles}）により，
表面張力と圧力勾配の離散位置を一致させる．
この経路は寄生流れを抑える原理を説明するうえで重要だが，
拡散界面 CSF のモデル誤差そのものは残る．

第二は，高密度比・低 $We$ を主対象とする圧力ジャンプ型の経路である．
ここでは Young--Laplace ジャンプを各相の PPE に直接入れ，
HFE により界面近傍の片側データを構成し，
欠陥補正で相内 Poisson 問題を解く．
表面張力は予測段階の体積力としてではなく，
表面エネルギーの仮想仕事から得る面上の毛管補正量として補正段階へ渡す．
静的平衡で消える成分と非平衡界面を動かす成分を分けるため，
圧力勾配値域への射影は診断と平衡判定に限り，
実際の速度補正では成分体積制約を取り除いた面上の毛管補正量を用いる．
この区別により，節点速度へ戻してから界面を移流する不一致だけでなく，
非平衡界面を動かす毛管成分を数値的な射影で消してしまう不整合も避ける．

\noindent\textbf{圧力ジャンプ型スキームの導出の見取り図．}
本稿の標準二相スキームは，曲率式を経験的に補正して得るのではなく，
次の有限次元の制約力問題として導く．
\begin{enumerate}[noitemsep,leftmargin=2em]
  \item 連続界面では，Young--Laplace 則を
        表面エネルギー $E_\Gamma=\sigma|\Gamma|$ の体積制約付き第一変分として読む
        （式~\eqref{eq:surface_energy_virtual_work}--\eqref{eq:surface_energy_constrained_equilibrium}）．
  \item 離散界面では，再初期化後の形ではなく，保存形 CLS が実際に運んだ
        再初期化前の輸送終点 $q_T$ を毛管仕事の幾何状態にする
        （式~\eqref{eq:capillary_transport_endpoint}）．
  \item その幾何状態の表面エネルギー変分を，CLS 面輸送の線形化
        $T_f(q_T)$ を通じて面速度空間へ引き戻し，
        圧力仕事と同じ面計量で Riesz 表現 $s_f$ と成分体積反力 $B_h$ を定義する
        （式~\eqref{eq:capillary_surface_energy_riesz_cochain}--\eqref{eq:capillary_component_volume_reactions}）．
  \item 圧力反力は，PPE が使う発散 $D_f$ と同じ面計量に対する
        Riesz 随伴 $G_A$ として一意に決める．
        同じ発散を持つ別の面勾配代表は，PPE 解だけを見れば同等でも，
        圧力仕事としては同等ではない（§~\ref{sec:projection_derivation}）．
  \item 最後に，$s_f$ から圧力反力と成分体積反力で表せる静的反力を除き，
        残り $h_{\sigma,f}$ を非圧縮・体積保存の許容毛管駆動として得る
        （式~\eqref{eq:capillary_component_saddle}--\eqref{eq:capillary_corrected_cochain_and_drive}）．
\end{enumerate}
したがって本稿の毛管閉包は，真円・楕円などの形状名や
曲率 cap，damping，CFL 縮小によって選ぶ近似ではない．
同じ面空間上で「表面エネルギーの仕事」「圧力の仕事」
「成分体積制約の仕事」を同時に釣り合わせる離散変分原理から導く．

本稿で検証済みと主張する範囲は後続の検証章で確定する．
単相 NS，一流体 Balanced--Force 経路，高密度比静止液滴，
圧力ジャンプと離散平衡判定を含む有限時間二相計算，
および一様格子上の強変形 CLS 受動移流を検証対象とする．
一方で，非一様格子上の移動界面 CLS と，
長時間・完全連成 NS 移動界面の収束性は今後の検証対象として明示的に分ける．
また CCD の $\Ord{h^6}$ は内点ステンシルの設計次数であり，
境界条件や界面条件の処理による全体次数の律速は各章で別に評価する．
このように，設計次数，閉包条件，検証済み範囲を混同しないことが
本稿全体の読み方である．

以下 §~\ref{sec:why_hard}〜§~\ref{sec:roadmap} では，
なぜ単相流の方法がそのまま使えないか，どの困難がどの手法で解決されるかを順番に説明する．
全体アルゴリズムの俯瞰図は次節 §~\ref{sec:roadmap} 末尾の §~\ref{sec:algo_overview} に示す．

% ═══════════════════════════════════════════════════════════════════════════════
\subsection{単相流と気液二相流の本質的違い}
\label{sec:why_hard}
% ═══════════════════════════════════════════════════════════════════════════════

数値流体力学（CFD）を学んだ読者は，まず単相流（空気のみ，水のみ）の
Navier--Stokes 方程式を扱うのが一般的である．
気液二相流とは，気相（ガス）と液相（液体）が共存し，
その間に明確な界面を形成しながら運動する流体現象であり，
気泡塔・核沸騰・波浪・噴霧燃焼など，工学・自然界に広く現れる重要な現象である．
なぜ「2種類の流体が共存する」だけで，これほど困難が増すのか？
その答えは一言でいえば「界面の存在」に尽きる．

\begin{table}[htbp]
\centering
\caption{単相流と気液二相流の本質的違い}
\renewcommand{\arraystretch}{\PaperTableArrayStretch}
\begin{tabular}{lll}
\hline
 & \textbf{単相流} & \textbf{気液二相流} \\
\hline
物性値     & 空間・時間一定            & 界面を挟んで \textbf{不連続に変化} \\
密度       & 一様 $\rho$               & $\rho_l / \rho_g \approx 1000$（水/空気） \\
粘性       & 一様 $\mu$                & $\mu_l / \mu_g \approx 100$ \\
支配方程式 & 全領域で同一              & 各相で異なる方程式＋界面条件 \\
力の項     & 体積力のみ                & 体積力＋\textbf{界面の表面張力（線力）} \\
幾何情報   & 不要                      & 界面の位置・法線・\textbf{曲率} が必要 \\
\hline
\end{tabular}
\end{table}

% ─────────────────────────────────────────────────────────────────────────────
\subsubsection{問題の核心：界面上の不連続}
% ─────────────────────────────────────────────────────────────────────────────

図~\ref{fig:density_jump}に示すように，
気液界面を挟んで密度と粘性が急激に変化する．
不連続を鋭く扱う手法も存在するが，本稿の CLS/CSF 一流体経路では，
界面を格子内の有限厚さの遷移層として滑らかに近似する．
その近似の仕方によって計算の安定性・精度・質量保存性が大きく変わる．

\begin{figure}[htbp]
\centering
\begin{tikzpicture}[scale=1.0]
  %% 左側：理想（物理的）界面
  \begin{scope}[xshift=-45mm]
    \node[font=\small\bfseries, text=navy] at (0, 2.5) {（a）物理的界面（理想）};
    \fill[blue2!15]   (-1.5,-2) rectangle (0,2);
    \fill[orange2!15] (0,-2)    rectangle (1.5,2);
    \draw[very thick, red2] (0,-2) -- (0,2);
    \draw[thick, blue2]
      (-1.5,-1.5) -- (-0.05,-1.5)
      (0.05,-1.5) -- (1.5,-1.5) node[right, font=\scriptsize]{$\rho_g$};
    \draw[dashed, red2, very thick] (0,-2) -- (0,-1.5);
    \node[font=\scriptsize, blue2] at (-0.75,-1.2) {$\rho_l$};
    \draw[<->, thick, darkgray]
      (-1.5,-1.8) -- (0,-1.8) node[midway, below, font=\scriptsize]{液相 $\Omega_l$};
    \draw[<->, thick, darkgray]
      (0,-1.8) -- (1.5,-1.8) node[midway, below, font=\scriptsize]{気相 $\Omega_g$};
    \node[font=\scriptsize, red2]  at (0.4, 0)  {界面 $\Gamma$};
    \node[font=\scriptsize]        at (-0.75, 1.0) {$\rho_l \approx 1000\rho_g$};
    \node[font=\scriptsize]        at (-0.75, 0.5) {$\mu_l \approx 100\mu_g$};
  \end{scope}
  %% 右側：数値的界面（CLS 遷移層）
  \begin{scope}[xshift=45mm]
    \node[font=\small\bfseries, text=navy] at (0, 2.5) {（b）数値的界面（CLS 遷移層）};
    \fill[blue2!15]   (-1.5,-2) rectangle (-0.4,2);
    \fill[orange2!15] (0.4,-2)  rectangle (1.5,2);
    \fill[teal2!10]   (-0.4,-2) rectangle (0.4,2);
    \draw[very thick, teal2, domain=-1.5:1.5, samples=60]
      plot (\x, {-1.5 + 0.7*(1/(1+exp(-5*\x)))});
    \draw[<->, thick, red2]
      (-0.4, 1.5) -- (0.4, 1.5)
      node[midway, above, font=\scriptsize, red2]{中心遷移幅 $\approx 2\varepsilon$};
    \draw[dashed, gray] (-0.4,-2) -- (-0.4, 1.4);
    \draw[dashed, gray] (0.4,-2)  -- (0.4, 1.4);
    \node[font=\scriptsize, teal2]   at (0,    0.2) {$\psi = H_\varepsilon(-\phi)$};
    \node[font=\scriptsize, blue2]   at (-1.0,-0.5) {$\psi\approx 1$（液相）};
    \node[font=\scriptsize, orange2] at (1.0, -0.5) {$\psi\approx 0$（気相）};
  \end{scope}
\end{tikzpicture}
\caption{気液界面における密度・粘性のジャンプ}
\label{fig:density_jump}
\PaperCaptionNote{左は物理的不連続，右は CLS 遷移層による数値表現．CLS 変数 $\psi \in [0,1]$ は界面を滑らかな遷移層として表現する．本稿では $\varepsilon$ を界面幅パラメータと呼び，中心遷移幅は $\approx 2\varepsilon$，95\% 遷移帯はおよそ $|\phi|\le 3\varepsilon$（幅 $\approx 6\varepsilon$）である．\textbf{符号規約：}液相で $\psi \approx 1$，気相で $\psi \approx 0$（本稿全章共通）．これは物性補間式 $\rho(\psi) = \rho_g + (\rho_l - \rho_g)\psi$ と整合する：$\psi=1$ のとき $\rho = \rho_l$（液相密度），$\psi=0$ のとき $\rho = \rho_g$（気相密度）．変数 $\phi$（符号付き距離関数）と $\psi$（保存形 Level Set 変数）の正式な定義は第~\ref{sec:governing}章（§~\ref{sec:notation}）で述べる．}
\end{figure}

% ─────────────────────────────────────────────────────────────────────────────
\subsubsection{数値拡散と界面の「にじみ」（Smearing）：視覚的理解}
% ─────────────────────────────────────────────────────────────────────────────

単純な1次精度の上流差分で界面を移流した場合，
時間が経つにつれて界面が広がり（にじみ），
密度・粘性の遷移幅が拡大していく．

界面の鮮鋭さが失われると：
\begin{itemize}
  \item 曲率 $\kappa$ の計算精度が低下する
  \item 密度の勾配が弱まり，表面張力が正確に評価できなくなる
  \item 液相体積（質量）が徐々に減少する（質量損失）
\end{itemize}

本稿の FCCD 面ジェット移流（Flux-form CCD；§~\ref{sec:advection_motivation}）と
Ridge--Eikonal 再初期化（§~\ref{sec:reinit}）は，
このにじみを抑制するための中心的な手法である．
FCCD は面値 $(\psi u)_{f,i+1/2}$ を Hermite 補間で滑らかに構成して保存形有限体積法の望ましい差し引き構造を保ち，
界面を壊す不要な高波数モードを構造的に抑える一方で，
段階 D の Ridge--Eikonal 距離再構成と段階 F の $\phi$ 空間体積閉鎖により，
拡散 CLS 質量と閾値化相体積を分けて閉じながら，
崩れた界面プロファイルを $\psi = H_\varepsilon(-\phi)$ の形へ戻す（§~\ref{sec:reinit}）．

% ═══════════════════════════════════════════════════════════════════════════════
\subsection{界面に起因する 4 つの数値的困難}
\label{sec:challenges}
% ═══════════════════════════════════════════════════════════════════════════════

気液二相流を数値的に扱う際には，以下の4つの本質的困難が存在する：

\begin{enumerate}
  \item \textbf{界面の捕捉と質量保存}：
        密度・粘性が界面を挟んで不連続に変化する
        （$\rho_l/\rho_g \approx 1000$，$\mu_l/\mu_g \approx 100$）ため，
        界面位置を高精度かつ安定に追跡しなければならない．
        標準的な Level Set 移流方程式 $\partial\phi/\partial t + \bu\cdot\bnabla\phi = 0$~\cite{OsherSethian1988} は
        非保存形であるため，数値積分誤差が体積損失として蓄積する~\cite{Sussman1999}．
        さらに移流スキームの数値拡散により界面プロファイルが「にじみ」，
        曲率計算精度が低下する．

  \item \textbf{表面張力の正確な評価}：
        表面張力 $\sigma\kappa_{lg}\hat{\bm n}_{lg}$ は本来「界面上の面力（2次元面力）」であり，
        体積方程式（Navier--Stokes）に直接組み込めない．
        曲率 $\kappa$ の計算には二階偏微分が必要となり，数値誤差が増幅されやすい．
        デルタ関数的な集中力をどう離散化するかが精度を決定する．

  \item \textbf{非圧縮条件の高精度維持}：
        速度場の発散をゼロに保つ圧力場を毎ステップ正確に求解する必要がある．
        単相流では均一係数のラプラシアン $\bnabla^2 p = \cdots$ であるが，
        二相流では可変係数 $\bnabla\cdot(1/\rho\,\bnabla p) = \cdots$ となり，
        界面をまたぐ密度比（$\rho_l/\rho_g \approx 1000$）が係数の急峻な変化を生み，
        反復ソルバーの収束が著しく悪化する．
        さらにコロケート格子では圧力と速度の数値的デカップリング
        （チェッカーボード不安定）が追加の困難をもたらす．

  \item \textbf{密度ジャンプと寄生流れの安定処理}：
        静水圧，表面張力ジャンプ，圧力投影を同じ離散位置で評価しないと，
        静止すべき液滴の界面に沿って渦状の人工的速度場（寄生流れ，spurious currents）が発生する~\cite{Francois2006, Sussman2000}．
        拡散界面 CSF 経路では，圧力勾配項 $-\bnabla p$ と
        表面張力項 $\sigma\kappa\bnabla\psi$ の演算子不一致が主因となる．
        分相圧力ジャンプ経路では，界面ジャンプ，PPE 補正，射影後の面フラックスが
        同じ面上の離散構造に閉じていないことが同等の不整合を生む．
        連続体レベルで力が釣り合っていても，離散レベルで残った不整合誤差が体積力として作用し，
        人工的な速度場を駆動する．
\end{enumerate}

4 つの困難はすべて次節の失敗例 1--4（§~\ref{sec:failure_modes}）のいずれかに発現する：
失敗例 1（寄生流れ）は\textbf{困難 2}（曲率の数値誤差）と\textbf{困難 4}（離散圧力--表面張力整合）が
共起する複合症状，失敗例 2（質量損失）と失敗例 4（にじみ）はいずれも\textbf{困難 1}
（界面捕捉）の異なる発現，失敗例 3（チェッカーボード）は\textbf{困難 3}
（非圧縮条件と圧力--速度結合）の典型である．本稿の対処方針は同節でまとめて提示する．

% ═══════════════════════════════════════════════════════════════════════════════
\subsection{二相流数値計算の典型的失敗例}
\label{sec:failure_modes}
% ═══════════════════════════════════════════════════════════════════════════════

上記の困難を軽視すると，計算は「動く」が物理的に意味のない結果を出力する．
実装者が遭遇しやすい4つの失敗例を示す．
以下の失敗例 1--4 はそれぞれ §~\ref{sec:challenges} の
困難 2 + 困難 4（曲率精度 \& BF 整合の共起＝寄生流れ）・
困難 1（質量損失）・困難 3（チェッカーボード）・困難 1（にじみ）に対応し，
各例の末尾に本稿の対処を明示する．

\paragraph{失敗例 1：寄生流れ（Spurious Currents）}
\subparagraph{症状．}静止しているはずの液滴が，界面に沿って渦状の速度場を持つ．
速度の大きさは表面張力に比例して大きくなる．

\subparagraph{原因．}曲率 $\kappa$ の計算誤差と圧力--表面張力の離散位置不一致が，
表面張力項 $\sigma\kappa\bnabla\psi$ または圧力ジャンプ条件
$j_{gl}=p_g-p_l=-\sigma\kappa_{lg}$ の残差として現れる．
この残差は圧力場と速度場を通じてフィードバックループを形成し，
数値的な「流れ」として固定化する．

\subparagraph{本稿の対処．}
CCD/FCCD による滑らかな領域の高次幾何評価に加え，
HFE による界面近傍データ延長，向き付き圧力ジャンプ分相 PPE，射影後の面速度，
履歴圧力のアフィンジャンプ面加速度を同じ面中心演算子族で閉じる．
一括 CSF/BF 一流体経路では $\bnabla p$ と $\kappa\bnabla\psi$ の演算子整合により
寄生流れを CSF モデル誤差の水準まで抑制し，
分相圧力ジャンプ経路では界面ジャンプ・補正・界面移流の面位置不一致を直接取り除く
（付録~\ref{sec:ppe_discretization_choice}）．
どこまでを検証済みとみなすかは第~\ref{sec:verification}章でまとめ，
非一様格子上の移動界面 CLS および長時間・完全連成 NS 移動界面は
§~\ref{sec:future_work} の今後の検証対象として残す．

\paragraph{失敗例 2：質量損失（Mass Loss）}
\subparagraph{症状．}時間が経つにつれて液滴が小さくなり，最終的に消える．

\subparagraph{原因．}標準的な Level Set 法では，移流方程式
$\partial\phi/\partial t + \bu\cdot\bnabla\phi=0$ が
保存形ではないため，数値積分誤差が体積損失として蓄積する．
再初期化の過程で $\phi=0$ 等値面が微小移動することも体積を変化させる
（定量化は §~\ref{sec:why_cls}）．

\subparagraph{本稿の対処．}Conservative Level Set 法（CLS）の保存形移流に加え，
段階 D の Ridge--Eikonal 距離再構成と段階 F の $\phi$ 空間体積閉鎖により，
拡散 CLS 質量 $\sum\psi_i\Delta V_i$ と閾値化相体積
$|\{\psi\ge 1/2\}|$ を分けて閉じ，
幾何プロファイル修復を物理移流の毛管仕事と混同しない．

\paragraph{失敗例 3：圧力--速度のデカップリング（Checkerboard 不安定）}
\subparagraph{症状．}圧力場が空間的に激しく振動し（チェッカーボードパターン），
速度場が非物理的な振動を示す．

\subparagraph{原因．}コロケート格子において，中心差分による圧力勾配の評価では
隣り合う格子点の圧力が連立されないため（1つ飛びの点しか参照しない），
高周波の圧力振動が速度場に伝播しない「幽霊モード」が存在する
（コロケート格子上のこの課題とその運動量補間による古典的対処は~\cite{RhieChow1983}）．
具体的に周期 $2\Delta x$ のチェッカーボードモード $p_j = P_0(-1)^j$ に対して，
中心差分は
$(p_{j+1}-p_{j-1})/(2\Delta x) = P_0\bigl[(-1)^{j+1}-(-1)^{j-1}\bigr]/(2\Delta x) = 0$
となりゼロ勾配を返す．より詳細な固有値解析は付録~\ref{app:checkerboard_mode} を参照．

\subparagraph{本稿の対処．} Balanced--Force 演算子整合（圧力勾配と CSF/界面力に同一の離散微分演算子を共有）と
FCCD 面ジェット部分系により，$2\Delta x$ チェッカーボードモードをスキーム内部で抑制する
（第~\ref{sec:collocate}章§~\ref{sec:checkerboard_solution} 参照）．

\paragraph{失敗例 4：界面のにじみ（Interface Smearing）}
\subparagraph{症状．}時間とともに界面が広がり，密度・粘性の遷移領域が厚くなっていく．

\subparagraph{原因．}移流スキームの数値拡散（numerical diffusion）により，
界面指標 $\psi$ のプロファイルが徐々にぼやける~\cite{OlssonKreiss2005}．
特に低次精度スキーム（1次風上差分など）では顕著．

\subparagraph{本稿の対処．}FCCD 面ジェット移流（§~\ref{sec:advection_motivation}）と
Ridge--Eikonal 再初期化（§~\ref{sec:reinit}）により，界面の鮮鋭さを維持する．

% ─────────────────────────────────────────────────────────────────────────────
\paragraph{本章と次章の繋がり．}
% ─────────────────────────────────────────────────────────────────────────────

以上 4 つの失敗例は，いずれも連続論では成立する保存則・整合条件が
離散化で崩れることに起因する．これらを統一的に扱うための基盤として，
第~\ref{sec:governing}章では支配方程式（Navier--Stokes と界面移流）を
One-Fluid 形で与え，応力ジャンプ条件の CSF 体積力への吸収機構を示す．
第~\ref{sec:levelset}章では CLS 法による $\phi \to \psi = H_\varepsilon(-\phi)$
平滑化変数の導入と保存形移流の定式化を述べ，以降の離散化設計に接続する．
